% Options for packages loaded elsewhere
\PassOptionsToPackage{unicode}{hyperref}
\PassOptionsToPackage{hyphens}{url}
\documentclass[
  ignorenonframetext,
]{beamer}
\newif\ifbibliography
\usepackage{pgfpages}
\setbeamertemplate{caption}[numbered]
\setbeamertemplate{caption label separator}{: }
\setbeamercolor{caption name}{fg=normal text.fg}
\beamertemplatenavigationsymbolsempty
% remove section numbering
\setbeamertemplate{part page}{
  \centering
  \begin{beamercolorbox}[sep=16pt,center]{part title}
    \usebeamerfont{part title}\insertpart\par
  \end{beamercolorbox}
}
\setbeamertemplate{section page}{
  \centering
  \begin{beamercolorbox}[sep=12pt,center]{section title}
    \usebeamerfont{section title}\insertsection\par
  \end{beamercolorbox}
}
\setbeamertemplate{subsection page}{
  \centering
  \begin{beamercolorbox}[sep=8pt,center]{subsection title}
    \usebeamerfont{subsection title}\insertsubsection\par
  \end{beamercolorbox}
}
% Prevent slide breaks in the middle of a paragraph
\widowpenalties 1 10000
\raggedbottom
\AtBeginPart{
  \frame{\partpage}
}
\AtBeginSection{
  \ifbibliography
  \else
    \frame{\sectionpage}
  \fi
}
\AtBeginSubsection{
  \frame{\subsectionpage}
}
\usepackage{iftex}
\ifPDFTeX
  \usepackage[T1]{fontenc}
  \usepackage[utf8]{inputenc}
  \usepackage{textcomp} % provide euro and other symbols
\else % if luatex or xetex
  \usepackage{unicode-math} % this also loads fontspec
  \defaultfontfeatures{Scale=MatchLowercase}
  \defaultfontfeatures[\rmfamily]{Ligatures=TeX,Scale=1}
\fi
\usepackage{lmodern}
\ifPDFTeX\else
  % xetex/luatex font selection
\fi
% Use upquote if available, for straight quotes in verbatim environments
\IfFileExists{upquote.sty}{\usepackage{upquote}}{}
\IfFileExists{microtype.sty}{% use microtype if available
  \usepackage[]{microtype}
  \UseMicrotypeSet[protrusion]{basicmath} % disable protrusion for tt fonts
}{}
\makeatletter
\@ifundefined{KOMAClassName}{% if non-KOMA class
  \IfFileExists{parskip.sty}{%
    \usepackage{parskip}
  }{% else
    \setlength{\parindent}{0pt}
    \setlength{\parskip}{6pt plus 2pt minus 1pt}}
}{% if KOMA class
  \KOMAoptions{parskip=half}}
\makeatother
\usepackage{longtable,booktabs,array}
\newcounter{none} % for unnumbered tables
\usepackage{calc} % for calculating minipage widths
\usepackage{caption}
% Make caption package work with longtable
\makeatletter
\def\fnum@table{\tablename~\thetable}
\makeatother
\setlength{\emergencystretch}{3em} % prevent overfull lines
\providecommand{\tightlist}{%
  \setlength{\itemsep}{0pt}\setlength{\parskip}{0pt}}
% Auto-generated by infrastructure/rendering/_pdf_latex_helpers.py::extract_math_font_preamble.
% Minimal header passed to Pandoc via -H so Beamer slides pick up the same math font
% as the combined-PDF path. Adjust the manuscript's preamble.md to change the font globally.
\usepackage{unicode-math}
\setmathfont{latinmodern-math.otf}

% Slide layout defaults for warning-clean scientific decks.
\usepackage{etoolbox}
\IfFileExists{xurl.sty}{\usepackage{xurl}}{}
\IfFileExists{seqsplit.sty}{\usepackage{seqsplit}}{\newcommand{\seqsplit}[1]{#1}}
\protected\def\breakseq#1{\seqsplit{#1}}
\protected\def\breaktt#1{\begingroup\ttfamily\seqsplit{#1}\endgroup}
\setlength{\emergencystretch}{6em}
\tolerance=5000
\hbadness=10000
\hfuzz=1pt
\setlength{\tabcolsep}{2pt}
\AtBeginEnvironment{longtable}{\tiny\renewcommand{\arraystretch}{0.86}\setlength{\tabcolsep}{1pt}}
\AtBeginEnvironment{tabular}{\tiny\renewcommand{\arraystretch}{0.86}\setlength{\tabcolsep}{1pt}}

% Natbib and cross-reference fallbacks — slides don't load natbib
% or cleveref, but combined-PDF manuscript prose may emit these
% commands. The fallback renders citations as a bracketed key list
% and cross-references as detokenized labels so slides don't fail on
% undefined control sequences or raw underscores. \providecommand is
% a no-op if packages load later.
\providecommand{\citep}[1]{[#1]}
\providecommand{\citet}[1]{#1}
\providecommand{\citealp}[1]{#1}
\providecommand{\citeauthor}[1]{#1}
\providecommand{\citeyear}[1]{#1}
\providecommand{\cref}[1]{\texttt{\detokenize{#1}}}
\providecommand{\Cref}[1]{\texttt{\detokenize{#1}}}

\newtheorem{proposition}{Proposition}
\newtheorem{hypothesis}{Hypothesis}
\usepackage{bookmark}
\IfFileExists{xurl.sty}{\usepackage{xurl}}{} % add URL line breaks if available
\urlstyle{same}
% fallback for those not using the hyperref driver hyperxmp:
\makeatletter
\@ifundefined{xmpquote}{\newcommand{\xmpquote}[1]{#1}}{}
\makeatother
\hypersetup{
  hidelinks,
  pdfcreator={LaTeX via pandoc}}

\author{\texorpdfstring{}{}}
\date{}

\begin{document}

\section{Experimental Setup}\label{sec:experimental_setup}

\begin{frame}[fragile,allowframebreaks]{Experimental Setup}
This section details the complete experimental configuration used to
generate the results in this study. Every parameter value below is
injected programmatically from \texttt{config.yaml} and the analysis
pipeline --- no value is hardcoded in the manuscript source.
\end{frame}

\begin{frame}[allowframebreaks]{Problem Definition}
\protect\phantomsection\label{problem-definition}
The optimization target is the quadratic objective function defined in
\breakseq{[@eq:quadratic\_objective]}:

\begin{equation}
\label{eq:quadratic_objective}
f(x) = \frac{1}{2} x^T A x - b^T x
\end{equation}

with \(A = [(1.0,)]\) and \(b = [1.0]\), yielding the analytical optimum
at \(x^* = 1.0\) with \(f(x^*) = -0.5\).
\end{frame}

\begin{frame}[allowframebreaks]{Parameter Space}
\protect\phantomsection\label{parameter-space}
The experiment systematically varies the gradient descent step size
across 6 values:

\begin{itemize}
\tightlist
\item
  \(\alpha = 0.01\) (conservative)
\item
  \(\alpha = 0.1\) (conservative)
\item
  \(\alpha = 0.5\) (near-optimal)
\item
  \(\alpha = 1.0\) (near-optimal)
\item
  \(\alpha = 1.5\) (aggressive)
\item
  \(\alpha = 2.5\) (divergent (expected unstable for H = I))
\end{itemize}

All runs start from the initial point \(x_0 = 0.0\) and use a
convergence tolerance of \(\|{\nabla f}\| < 10^{-8}\) with a maximum
iteration limit of \(N_{\max} = 1000\).
\end{frame}

\begin{frame}[fragile,allowframebreaks]{Numerical Stability Grid}
\protect\phantomsection\label{numerical-stability-grid}
To validate robustness, the optimizer is exercised across a grid of 8
starting points and 6 step sizes, producing 48 total evaluations. This
comprehensive sweep confirms that convergence is not an artifact of a
narrow parameter choice. The stability metric is computed for the
\breaktt{quadratic\_function} objective via
\breaktt{infrastructure.scientific.stability.check\_numerical\_stability()}.
\end{frame}

\begin{frame}[fragile,allowframebreaks]{Dimensional Scaling}
\protect\phantomsection\label{dimensional-scaling}
Performance benchmarking spans problem dimensions
\(d \in \{1, 2, 5, 10, 20, 50\}\), from the scalar case (\(d = 1\)) to
moderate dimensionality (\(d = 50\)), using identity-Hessian quadratics
to isolate algorithmic scaling from problem conditioning effects.
Representative single-call execution time from the last benchmark run:
\textbf{1.9 μs} (recorded in
\breaktt{output/reports/performance\_benchmark.json}).
\end{frame}

\begin{frame}[allowframebreaks]{Computational Environment}
\protect\phantomsection\label{computational-environment}
\begin{itemize}
\tightlist
\item
  \textbf{Python}: 3.12.13
\item
  \textbf{NumPy}: 2.4.2
\item
  \textbf{Platform}: Darwin arm64
\item
  \textbf{Generated}: 2026-07-10T14:33:59Z
\end{itemize}
\end{frame}

\begin{frame}[fragile,allowframebreaks]{Pipeline ordering}
\protect\phantomsection\label{pipeline-ordering}
Typical \breaktt{template\_code\_project} analysis order (see
\breaktt{scripts/pipeline/stage\_02\_analysis.py} discovery) is:

\begin{enumerate}
\tightlist
\item
  \breaktt{optimization\_analysis.py} --- writes
  \breaktt{output/data/optimization\_results.csv},
  \breaktt{output/figures/*.png}, and JSON reports under
  \texttt{output/reports/}.
\item
  \breaktt{z\_generate\_manuscript\_variables.py} --- reads the CSV and
  YAML, emits \breaktt{output/data/manuscript\_variables.json}, and
  writes substituted copies to \breaktt{output/manuscript/} for
  rendering.
\item
  \breaktt{generate\_api\_docs.py} --- refreshes API markdown consumed by
  documentation targets.
\end{enumerate}

PDF compilation then reads from \breaktt{output/manuscript/} so that
figure paths and numeric tables match the analysis that just completed.
\end{frame}

\begin{frame}[fragile,allowframebreaks]{Relation to figures}
\protect\phantomsection\label{relation-to-figures}
{\def\LTcaptype{none} % do not increment counter
\begin{longtable}[]{@{}
  >{\raggedright\arraybackslash}p{(\linewidth - 2\tabcolsep) * \real{0.5263}}
  >{\raggedright\arraybackslash}p{(\linewidth - 2\tabcolsep) * \real{0.4737}}@{}}
\toprule\noalign{}
\begin{minipage}[b]{\linewidth}\raggedright
Figure (\breakseq{[@sec:results]})
\end{minipage} & \begin{minipage}[b]{\linewidth}\raggedright
Primary inputs
\end{minipage} \\
\midrule\noalign{}
\endhead
Convergence trajectories & \breaktt{experiment.step\_sizes},
\texttt{initial\_point}, \texttt{tolerance}, \texttt{max\_iterations} \\
Step-size sensitivity & Dense \(\alpha\) sweep internal to
\breaktt{generate\_step\_size\_sensitivity\_plot()} \\
Convergence rate & Same trajectories as above; tolerance line uses
\breaktt{convergence\_tolerance} \\
Complexity quad panel & One bar chart per row of
\breaktt{optimization\_results.csv} \\
Stability heatmap & \breaktt{stability\_starting\_points} \(\times\)
\breaktt{stability\_step\_sizes} \\
Dimensional benchmark & \breaktt{benchmark\_dimensions}, fixed
\(\alpha=0.1\), internal tol \(10^{-10}\) \\
\bottomrule\noalign{}
\end{longtable}
}

This table is descriptive documentation only; it is not executed as code
during the build.
\end{frame}

\end{document}
